%% kapa: exam=1 round=33 year=2022 period=2 form=A subject=accounting subject_ko=회계학 seq=37 number=77
%% source: Q-Net 감정평가사 기출문제 (cst00309) artlSeq=5209997 file=raw/1008/5209997/2022년도 제33회 감정평가사 1차시험 2교시(감정평가 관계법규, 회계학).pdf
%% conversion: kordoc 4.24.51 -> tools/gen1.py (auto); review=proofread-sonnet; flags=none
\begin{question}{77}
다음은 (주)감평의 20×1년 상반기 종합예산을 작성하기 위한 자료의 일부이다. 4월의 원재료 구입예산액은?
\begin{framedtext}
○ 예산판매량 \par
- 3월: 2,000단위 \quad 4월: 2,500단위 \quad 5월: 2,400단위 \quad 6월: 2,700단위 \par
○ 재고정책 \par
- 제품: 다음 달 예산판매량의 10\%를 월말재고로 보유한다. \par
- 원재료: 다음 달 생산량에 소요되는 원재료의 5\%를 월말재고로 보유한다. \par
○ 제품 1단위를 생산하는데 원재료 2kg이 투입되며, kg당 구입단가는 ￦10이다. \par
\end{framedtext}
\begin{choices}
  \item ￦49,740
  \item ￦49,800
  \item ￦49,860
  \item ￦52,230
  \item ￦52,290
\end{choices}
\end{question}
\answer{1}
